\documentclass[10pt,a4paper]{amsart}
\usepackage[margin=19mm]{geometry}
\usepackage{amsmath,amssymb,mathtools}
\usepackage[scr=rsfs,cal=euler]{mathalfa}
\usepackage{xfrac}
\usepackage[unicode,hidelinks,bookmarks=false]{hyperref}
\newcommand{\ie}{i.e.\ }
\newcommand{\cf}{cf.\ }
\newcommand{\resp}{respectively\ }
\newcommand{\Subsection}{\subsection}
\providecommand{\nobreakdash}{\nobreak\hbox{-}}
\setlength{\emergencystretch}{2em}
\multlinegap0pt
\usepackage{bm}
\usepackage{bbm}
\usepackage{dsfont}
\usepackage{xspace}
\usepackage{cancel}
\usepackage{mathtools}
\usepackage{amsthm}
\makeatletter
\@ifundefined{theorem}{\newtheorem{theorem}{Theorem}}{}
\@ifundefined{lemma}{\newtheorem{lemma}[theorem]{Lemma}}{}
\makeatother
\title[Wall And Chamber Structure For A Special Biserial Algebra Co]{Wall And Chamber Structure For A Special Biserial Algebra Coming From Perverse Sheaves on $\mathbb{P}^n$}
\author{Alessio Cipriani, Martina Lanini}
\date{Source: arXiv:2308.00991v2 (CC BY 4.0)}
\begin{document}
\maketitle

\noindent\emph{Source and licence.} Wall And Chamber Structure For A Special Biserial Algebra Coming From Perverse Sheaves on $\mathbb{P}^n$. Version arXiv:2308.00991v2 (CC BY 4.0), distributed under the Creative Commons Attribution 4.0 International licence, as recorded in the arXiv licence metadata for this exact version and shown on the abstract page https://arxiv.org/abs/2308.00991v2. Full rights notice: licenses/arxiv-2308.00991-CC-BY-4.0.txt.\par\smallskip

\noindent\textit{Editorial scope.} Counted unit: the Lemma whose source label is \texttt{ lemma : parameterisation of strings} in the article (source0.tex lines 376--381, source label \texttt{ lemma : parameterisation of strings}), delivered together with the article's own complete original proof of it (source0.tex lines 382--405, one \texttt{proof} environment of 24 lines). No mathematics is written, repaired or re-derived: statement and proof are the article's own text line for line. Required context retained uncounted: none. Every definition, display and label of those lines is retained unchanged. Omitted from this package and not redistributed: 1--375, 406--1324. Nothing inside a retained excerpt is omitted or altered: every retained body is delivered exactly as the article's lines are written, so no omission receipt of any kind accompanies this package. Citations: the retained excerpts carry no citation command at all, so no bibliography entry of the article is transcribed and no reference is redistributed. Reference adaptation: none was needed and none was made; every reference inside the delivered unit resolves inside it, so no unresolved reference or citation is produced. Printed slips kept literal: no printed word, symbol, hypothesis or number of the counted statement or of its proof was altered. Layout and class: the article's own class is \texttt{amsart} with packages including amsthm, amssymb, latexsym, euscript, amsmath, amstext, amsgen, amsbsy, amsopn, amsfonts, hyperref, graphicx, bbm, tikz; the wrapper is the lane's pinned \texttt{amsart} editorial wrapper at 10pt on A4, which restates the theorem-like environments the retained text needs and adds the article's own macro definitions that the retained text needs (none), with extra wrapper packages (bm, bbm, dsfont, xspace, cancel, mathtools). The delivered numbering is the unit's own. Assets: no retained span contains a figure environment, a graphics-inclusion command, a PDF-page inclusion, a table or a TikZ picture; no image file is copied into this package, the package declares no asset dependency and no image file is redistributed with it. Licence: version arXiv:2308.00991v2 of this article is distributed under the Creative Commons Attribution 4.0 International licence, as recorded in the arXiv licence metadata for this exact version and shown on the abstract page https://arxiv.org/abs/2308.00991v2. A full rights notice is delivered at licenses/arxiv-2308.00991-CC-BY-4.0.txt. Characters: every retained excerpt is pure ASCII, so the delivered engine, XeLaTeX with its default Latin Modern text font, reports no missing character.
\begin{lemma}\label{ lemma : parameterisation of strings}%Let $r\geq 2$. 
The set of $*$-classes of alternating walk strings on $(Q(n), I(n))$ is in bijection with the following set
    \[
    \mathscr{S}=\left\{(a,b,\eta)\mid 0\leq a<b\leq n, \ \eta\in \{\pm1\} \right\}.
    \]
\end{lemma}

\begin{proof}
    We denote by $\mathscr{X}$ the set of $*$-classes of strings on $(Q(n), I(n))$ of length $>0$. 
    We will define a bijective function  $\varphi:\mathscr{X}\rightarrow \mathscr{S}$.

    Let $w=\gamma_m\ldots \gamma_1$ be a string of length $>0$. Then $f_w$ is strictly monotone and it is increasing if and only if $f_{w^*}$ is decreasing. Since we only care about $*$-classes we can assume that $f_w$ is strictly increasing. We define 
    \[
    \varphi(w)=(f_w(0), f_w(m), \eta), \hbox{ with }\eta=\begin{cases}
        1&\hbox{ if }\gamma_1\in Q(n)_1,\\
        -1&\hbox{ if }\gamma_1\in Q(n)_1^*.
    \end{cases}
    \]
    Viceversa, if $(a,b,\eta)\in\mathscr{S}$ we set
    \[
 \psi((a,b,\eta))=\begin{cases}
     \alpha_{b-1}^*\beta_{b-2}\ldots\alpha_{a+1}^*\beta_a &\hbox{ if }\eta=1, \hbox{ and }a\equiv_2 b\\
     \beta_{b-1}\alpha_{b-2}^*\ldots\alpha_{a+1}^*\beta_a &\hbox{ if }\eta=1, \hbox{ and }a\not\equiv_2 b\\
     \beta_{b-1}\alpha_{b-2}^*\ldots\beta_{a+1}\alpha_a^* &\hbox{ if }\eta=-1, \hbox{ and }a\equiv_2 b\\
     \alpha_{b-1}^*\beta_{b-2}\ldots\beta_{a+1}\alpha_{a}^* &\hbox{ if }\eta=-1, \hbox{ and }a\not\equiv_2 b,
 \end{cases}
    \]
    where $a\equiv_2 b$, resp. $a\not\equiv_2 b$, indicates that $a$ and $b$ have, resp. do not have, the same parity.
    
  Since $t(\beta_i)=i+1=t(\alpha_{i+1})=s(\alpha_{i+1}^*)$ and $t(\alpha_j^*)=s(\alpha_j)=s(\beta_{j+1})$ then $(a,b,\eta)$ is certainly a reduced walk. It is a string since it does not contain any path of length $\geq 2$ or any monomial in $I(n)$ or summand of a binomial generating $I(n)$. Note that $s(\psi((a,b,\eta)))=a$ and $t(\psi((a,b,\eta)))=b$. At this point is an easy check to see that $\varphi$ and $\psi$ are mutual inverses.
\end{proof}

\end{document}
